\documentclass[a4paper,12pt]{article}
\usepackage[utf8]{inputenc}
\usepackage[russian]{babel}
\usepackage{amsmath, amssymb}

\begin{document}

\section{Факторизация многочлена $x^7 - 1$ над $\mathbb{F}_2$}

Рассмотрим разложение многочлена $x^7 - 1$ в поле $\mathbb{F}_2$. Поскольку в этом поле выполняется равенство $-1 = 1$, имеем:
\[
x^7 - 1 = x^7 + 1.
\]

Данный многочлен представим в виде произведения неприводимых множителей:
\[
x^7 + 1 = (x + 1)(x^3 + x + 1)(x^3 + x^2 + 1).
\]

Корректность разложения может быть проверена, например, с помощью символьных вычислений (например, в системе \texttt{sympy}).

\section{Получение порождающих многочленов}

Порождающие многочлены циклических кодов длины 7 соответствуют всем делителям $x^7 + 1$, за исключением тривиальных (единицы и самого многочлена).

Таким образом, получаем следующий набор:
\begin{align*}
g_1(x) &= x + 1, \\
g_2(x) &= x^3 + x + 1, \\
g_3(x) &= x^3 + x^2 + 1, \\
g_4(x) &= (x + 1)(x^3 + x + 1) = x^4 + x^3 + x^2 + 1, \\
g_5(x) &= (x + 1)(x^3 + x^2 + 1) = x^4 + x^2 + x + 1, \\
g_6(x) &= (x^3 + x + 1)(x^3 + x^2 + 1) \\
       &= x^6 + x^5 + x^4 + x^3 + x^2 + x + 1.
\end{align*}

Всего получено шесть различных нетривиальных порождающих многочленов.

\section{Параметры соответствующих кодов}

Для каждого из найденных многочленов можно определить параметры циклического кода длины $n = 7$.

Размерность вычисляется по формуле $k = 7 - \deg g(x)$, а минимальное расстояние определялось перебором всех кодовых слов.

\begin{center}
\begin{tabular}{|c|c|c|c|l|}
\hline
$g(x)$ & $\deg g$ & $k$ & $d_{\min}$ & Тип кода \\
\hline
$g_1$ & 1 & 6 & 2 & код чётности \\
$g_2$ & 3 & 4 & 3 & код Хэмминга (7,4) \\
$g_3$ & 3 & 4 & 3 & эквивалентный код Хэмминга \\
$g_4$ & 4 & 3 & 4 & симплекс-код \\
$g_5$ & 4 & 3 & 4 & эквивалентный симплекс-код \\
$g_6$ & 6 & 1 & 7 & код повторения \\
\hline
\end{tabular}
\end{center}

\section{Интерпретация результатов}

Полученные коды представляют собой все возможные нетривиальные циклические двоичные коды длины 7.

Их можно интерпретировать следующим образом:

\begin{itemize}
    \item Многочлен $g_1(x)$ задаёт код, в котором все слова имеют чётный вес.
    
    \item Многочлены $g_2(x)$ и $g_3(x)$ определяют код Хэмминга $(7,4,3)$. Несмотря на различие в записи, соответствующие коды эквивалентны.
    
    \item Многочлены $g_4(x)$ и $g_5(x)$ задают симплекс-коды размерности 3. Эти коды являются двойственными к кодам Хэмминга.
    
    \item Многочлен $g_6(x)$ соответствует коду повторения, содержащему всего два кодовых слова.
\end{itemize}

Таким образом, перечисленные конструкции полностью описывают все циклические коды длины 7 над полем $\mathbb{F}_2$.

\end{document}